\section{Preliminaries}
\label{sec:preliminaries}
In this section, we provide a preliminary definition of terms which will be
used throughout the rest of this paper for use in the definitions of
the algebraic mapping operators.

We assume the following pairwise disjoint countably infinite sets: $\symAttrSet$
(attribute names), $\symAllIRIs$ (IRIs), $\symAllBNodes$ (blank nodes),
$\symAllLiterals$ (literals), and $\symFragSet$ (\fragment{} names).
We define a set of terms $\symTermSet$ as $\symTermSet = \symAllIRIs \cup
	\symAllBNodes \cup \symAllLiterals$.
Furthermore, we also define a special value $\error$ with 
$\error \notin (\symTermSet \cup \symAttrSet \cup \symFragSet)$. 


The \emph{schema of a mapping relation} is a tuple $\schemaTuple$, where $A
\subset \symAttrSet$ is a finite, non-empty set of attributes and $\fragAttr
\notin A$. 
An \emph{instance} for a schema of mapping relation, $\schemaTuple$, is a
multiset of so-called \emph{mapping tuples} where each such mapping tuple,
$\mappingTuple$, is a partial function 
$\mappingTuple: \symAttrSet \rightarrow \symFragSet \cup \symTermSet \cup
\{\error\} $ such that $\fctMappingTuple{\fragAttr} \in \symFragSet$ and 
$\forall a \in A. \fctMappingTuple{a} \in \symTermSet \cup \{\error\}$. 
A \emph{mapping relation} is a pair $\mappingRel$, where $\schema$ is a schema
of a mapping relation and $\mappingInst$ is an instance of $\schema$. 


